\section{语义与查询表示为何可能有效}
\noindent\textbf{本节结论与推导路径。}
上一节把困难定位为有限数据下的任务相关误差，而非笼统的像素信息不足。
由已有RGB确定的语义图不会增加原始观测信息，但外部分割监督与几何查询监督
可以提供任务相关的归纳偏置，在有限数据下减少表示丢失的信息和实际动作读出误差。
收益要求所选几何关系与控制有关、估计足够可靠，且策略能够利用该表示，
新增编码与监督的代价没有抵消这些收益。
本节先用条件期望将平方动作风险分解为不可约风险、表示损失与读出误差，
再由几何的近似充分性及Lipschitz条件，推导几何估计误差对动作风险的条件性上界，
并在同一任务敏感性加权分布下将其连接回上一节的任务界，
最后分析伪标签与不可见关系如何限制该结论。
因此，实验须区分几何可读出、动作实际使用该表示，以及闭环任务表现改善这三个层次。

\subsection{增加原始观测与加入语义先验不同}
在固定评价分布 $\mu$ 上，令 $Y$ 为相同目标执行时刻的归一化示范动作，
所有变量二阶可积，所有模型参数视为训练后固定。
若 $X$ 为已有输入，$I$ 为新增视角或更新的图像，则理想平方预测风险满足
\begin{equation}
 R^*(X)-R^*(X,I)
 =\E\norm{\E[Y\mid X,I]-\E[Y\mid X]}^2\ge0.
 \label{eq:newinfo}
\end{equation}
额外信息改变条件动作均值时才严格改善；实际学习误差可能抵消该收益。

若分割图 $S=g(I)$ 由已有RGB确定，固定分割器后
\begin{equation}
 \sigma(I,S)=\sigma(I),\qquad R^*(I,S)=R^*(I).
 \label{eq:semanticinfo}
\end{equation}
所以语义不是对原始RGB增加新的测量信息。
外部分割标签及预训练引入了额外的\emph{学习先验}，
将对象身份与空间范围直接编码，可能使有限数据模型更易学习合适映射。
这也说明只比示范条数不足以衡量公平性：额外标签、人工审核和预训练均应披露。

令 $\mathcal O$ 为共同允许的输入信息，$Z$ 为实际表示，
$m_{\mathcal O}=\E[Y\mid\mathcal O]$，$m_Z=\E[Y\mid Z]$。
条件期望正交性给出
\begin{align}
 R(p_Z)&=R(m_{\mathcal O})+a_Z+e_Z,\label{eq:decomp}\\
 a_Z&=\E\norm{m_{\mathcal O}-m_Z}^2,\quad
 e_Z=\E\norm{p_Z-m_Z}^2.\nonumber
\end{align}
语义的收益必须落在表示信息损失 $a_Z$ 和实际读出误差 $e_Z$ 的净下降上。
增加token与共享编码器也可能增加优化困难，不能只看输入更直观就预设收益。

\subsection{归纳偏置能处理哪类覆盖缺失}
将控制相关变量记为 $(G,C)$，无关外观记为 $N$。
若不同环境均满足
\begin{equation}
 \E[Y\mid G,C,N]=f(G,C),
 \label{eq:invariance}
\end{equation}
且变化只发生于 $N$，则同一几何决策函数无需重新学习每种纹理。
例如目标和工具关系不变但背景颜色改变，可靠语义可以减少对外观偶然相关性的依赖。
这要求语义估计器本身对该变化也可靠，且相关几何和控制条件已有覆盖。
它不解决从未见过的工具接触模式、物体动力学、目标完全消失后的任意恢复。

式\eqref{eq:invariance}是需验证的假设，不是分割标签自动获得的因果结论。
历史中上一动作、夹爪状态等也可能成为错误捷径\cite{dehaan2019}。
共享RGB路径并未被强制成为不变表示；应通过外观干预和几何保持对照检验，
不能仅用分割IoU解释策略的分布外表现。

\subsection{几何监督QToken的条件性作用}
令 $G$ 为任务几何，$C$ 为必要控制上下文。
假设存在关于几何为 $L_G$-Lipschitz的 $f$，使
\begin{equation}
 \E\norm{m_{\mathcal O}-f(G,C)}^2\le\epsilon_{\rm suf}.
\end{equation}
$\epsilon_{\rm suf}$ 容纳几何之外的姿态、接触、历史等因素。
从查询表示 $Z_Q$ 可读出 $\widehat G=d(Z_Q)$，并可访问 $C$，
设 $\epsilon_G=\E\norm{\widehat G-G}^2$。
\begin{proposition}[几何表示的误差界]
实际动作读出误差记为 $e_Q$，则
\begin{equation}
 R(p_Q)-R(m_{\mathcal O})
 \le(\sqrt{\epsilon_{\rm suf}}+L_G\sqrt{\epsilon_G})^2+e_Q.
 \label{eq:geometry}
\end{equation}
\end{proposition}
证明利用 $f(\widehat G,C)$ 作为 $Z_Q$ 的候选动作读出，见附录。
若基线的超额风险 $a_0+e_0$ 大于右侧，则有严格平方风险改善。
但基线表示信息损失的下界通常难以估计；一个较差的探针误差不是其下界。
该命题给出可能的充分条件，不证明加入辅助监督后必然满足这些条件。

语义图给出对象区域；几何监督query进一步要求指定关系能够被编码和读出。
具体监督关系由任务选择，而不要求一种固定的几何量或完整阶段标签。
辅助loss给query及其关联视觉特征提供梯度，
但不指定attention map，也不保证动作decoder实际利用了它。
因此需同时检查几何精度、query消融后的动作变化和闭环成功率。
不能把当前QToken写成已经具备阶段切换或停止判定的模型。

\subsection{把几何表示界接回任务性能界}
仅证明平均几何误差变小，尚未回答上一节提出的关键状态问题。
令 $Z_w=\E_{\bar d_E}w>0$，定义归一化的敏感性加权分布
$d\mu_w=(w/Z_w)d\bar d_E$。
本节的投影分解与几何界可统一在 $\mu_w$ 下重新推导，
但条件均值及所有误差也必须按该分布重新定义，不能只给最终公式乘一个权重。
以角标 $w$ 表示这些量，记
\begin{equation}
 \mathcal B_w=R_w^*+
 (\sqrt{\epsilon_{{\rm suf},w}}+L_G\sqrt{\epsilon_{G,w}})^2+e_{Q,w}.
 \label{eq:weightedrepresentation}
\end{equation}
若监督目标确为前节参考专家在同一决策状态、同一执行时刻的动作，
则 $\epsilon_{E,w}=Z_wR_w(p_Q)\le Z_w\mathcal B_w$，从而
\begin{equation}
 J(\pi_Q)\ge J(\pi_E)-T\sqrt{\bar w_{\pi_Q}C Z_w\mathcal B_w}.
 \label{eq:geometrytotask}
\end{equation}
$R_w^*$ 表示允许观察下仍然不可约的预测风险。
示范动作含模式歧义或目标时间错位时，不能不经校验就把监督风险代入此式。

这给出了层层连接的充分条件：任务敏感的控制关系在所选几何中近似充分，
其估计与读出误差足够小，且部署覆盖受控，才有更紧任务保证的可能。
严格改善几何表示的动作风险，要求加权几何界的超额部分
小于基线的加权超额风险；仅比较两个上界大小仍不是实际策略排序证明。
新增语义不被宣称自动降低 $C$、$w$ 或未覆盖质量 $\kappa$，
而是在这些条件成立的范围内改善表示相关风险。

因此，关键状态与QToken之间不是直接等号，而是
\emph{任务敏感性提出需要区分的关系，几何监督尝试让这些关系被保留，
动作读出再将其转为更准确的决策}。
监督了不相关关系，或关系虽准确却未被动作使用，都不能完成这一链条。

\subsection{伪标签、不确定性与不可见关系}
令 $\widetilde G$ 为伪标签。无偏假设不是必需的，但真值误差满足
\begin{equation}
 \epsilon_G\le
 2\E\norm{\widehat G-\widetilde G}^2+
 2\E\norm{\widetilde G-G}^2.
 \label{eq:pseudolabel}
\end{equation}
仅减小伪标签拟合误差不能保证第二项小；
稳定地把机械臂误分为目标可能同时具有高时间一致性和高系统偏差。
质量分数只能作为诊断和相对权重，须通过独立人工抽样验证。

对标量几何和质量权重 $W\ge0$，加权平方目标的最优读出为
\begin{equation}
 d_W(z)=\frac{\E[W\widetilde G\mid Z=z]}{\E[W\mid Z=z]},
 \quad \E[W\mid Z=z]>0.
 \label{eq:weighted}
\end{equation}
它不自动等于 $\E[G\mid Z=z]$。把困难状态长期降权可能降低其有效覆盖。
实际Smooth L1的最优解也不直接等于式\eqref{eq:weighted}。
对不可见关系，几何界只在有效域解释；目标不可见时面积零仍有效，
但目标与工具距离未定义，不能伪造零距离监督。
不可见域的策略风险必须单独报告，不能在整体成功率论证中删掉。

\subsection{表示改善之后仍然缺少什么}
即使上述几何误差和读出误差很小，固定观察下的 $R_w^*$ 仍可能为正：
任务相关动态不能总由当前静态几何辨别，
执行期间的新变化也不包含在生成原计划的旧观察中。
因此，单纯继续强化同一帧的语义监督不能消除全部残余误差。
下一节转向\emph{扩大可用的因果信息并学习其对当前计划的修正}：
更新观察响应执行中变化，历史帮助区分动态歧义。
这不是取代本节的表示学习，而是针对它没有解决的信息与时序缺口。
